(1) 规定逆时针为正电流，所围正法向为 \(+\hat z\)。磁通 \(\Phi=h\int_x^{x+w}(B_0+\beta s)\,ds=hw[B_0+\beta(x+w/2)]\)，因此动生电动势 \(\mathcal E=-d\Phi/dt=-h\beta wv\)。均匀背景 \(B_0\) 不贡献电动势，这揭示决定运动感应的是磁通的空间梯度而不是磁场的绝对大小。
(2) 忽略自感时 \(I=\mathcal E/R=-h\beta wv/R\)。两条竖直边所受洛伦兹合力的水平分量为 \(F_x=Ih[B(x+w)-B(x)]=Ih\beta w=-(h\beta w)^2v/R\)，永远与速度反向。消耗功率 \(I^2R=(h\beta w)^2v^2/R=-F_xv\)，恰好是动能被制动转化的热量。
(3) 回路仅受上述水平磁力，故 \(m\dot v=-\Gamma v\)，\(\Gamma=(h\beta w)^2/R>0\)。给定初速度 \(v_0\)，解为 \(v(t)=v_0 e^{-t/\tau}\)、\(\tau=m/\Gamma=mR/(h\beta w)^2\)；位移 \(x(t)-x(0)=v_0\tau(1-e^{-t/\tau})\)。当 \(\beta=0\) 时 \(\Gamma=0\)，回路匀速运动而不是发生虚假的制动。
(4) 自感存在时回路方程 \(L\dot I+RI=-h\beta wv\)，牛顿方程 \(m\dot v=Ih\beta w\)。将前者乘 \(I\)、后者乘 \(v\) 并相加，得 \(d(mv^2/2+LI^2/2)/dt=-RI^2\)。机械能损失的一部分可暂时储存在自感中，只有含两种能量的等式才是任意时间精确成立的守恒预算。
